% =============================================================
%  EEE 2188 – Electronics Sessional
%  Experiment 01: Familiarization of Equipment
%  Rajshahi University of Engineering & Technology (RUET)
% =============================================================

\documentclass[12pt, a4paper, twoside]{report}

% ---------------------------------------------------------------
% Core packages
% ---------------------------------------------------------------
\usepackage[T1]{fontenc}
\usepackage[utf8]{inputenc}
\usepackage{amsmath}
\usepackage{mathptmx}            % Times New Roman (text + math)
\usepackage{microtype}

% ---------------------------------------------------------------
% Page geometry
% ---------------------------------------------------------------
\usepackage[
  a4paper,
  top=2.5cm, bottom=2.5cm,
  inner=3.0cm, outer=2.5cm,
  headheight=15pt
]{geometry}

% ---------------------------------------------------------------
% Headers / Footers
% ---------------------------------------------------------------
\usepackage{fancyhdr}
\usepackage{lastpage}

\pagestyle{fancy}
\fancyhf{}

% EVEN pages  (left-hand page):  Course Code   Experiment No.
% ODD  pages  (right-hand page): Course Title  Current Section
\fancyhead[LE]{\small\itshape\CourseCode\ \ Experiment\ \ExperimentNo}
\fancyhead[RO]{\small\itshape\CourseTitle\ \ \textendash\ \ \rightmark}
\fancyfoot[LE,RO]{\small\thepage}
\renewcommand{\headrulewidth}{0.4pt}
\renewcommand{\sectionmark}[1]{\markright{#1}}

\fancypagestyle{plain}{%
  \fancyhf{}%
  \fancyfoot[LE,RO]{\small\thepage}%
  \renewcommand{\headrulewidth}{0pt}%
}

% ---------------------------------------------------------------
% Figures, tables, captions
% ---------------------------------------------------------------
\usepackage{graphicx}
\usepackage{xparse}
\let\OriginalIncludeGraphics\includegraphics
\RenewDocumentCommand{\includegraphics}{O{}m}{%
  \IfFileExists{#2}{\OriginalIncludeGraphics[#1]{#2}}{%
    \fbox{\parbox[c][5cm][c]{\linewidth}{\centering Missing image\\\detokenize{#2}}}%
  }%
}
\usepackage{float}
\usepackage[
  format=plain,
  labelfont=bf,
  font=small,
  justification=centering
]{caption}
\usepackage{subcaption}

% Chapter-based numbering: Fig. 1.1, Table 1.1 …
\usepackage{chngcntr}
\counterwithin{figure}{chapter}
\counterwithin{table}{chapter}

% ---------------------------------------------------------------
% Colours, boxes
% ---------------------------------------------------------------
\usepackage[dvipsnames]{xcolor}
\usepackage{tcolorbox}
\tcbuselibrary{skins, breakable}

% ---------------------------------------------------------------
% Hyperlinks
% ---------------------------------------------------------------
\usepackage[
  colorlinks=true,
  linkcolor=black,
  citecolor=black,
  urlcolor=blue,
  pdfborder={0 0 0}
]{hyperref}

% ---------------------------------------------------------------
% Bibliography – IEEE / ieeetr style
% ---------------------------------------------------------------
\usepackage[numbers,sort&compress]{natbib}
\bibliographystyle{ieeetr}

% ---------------------------------------------------------------
% Spacing / Paragraphs
% ---------------------------------------------------------------
\usepackage{setspace}
\setstretch{1.25}
\setlength{\parskip}{6pt}
\setlength{\parindent}{0pt}

% ---------------------------------------------------------------
% Enumerate / lists
% ---------------------------------------------------------------
\usepackage{enumitem}
\setlist[enumerate]{topsep=4pt, itemsep=2pt, leftmargin=1.8em}
\setlist[itemize]{topsep=4pt, itemsep=2pt, leftmargin=1.8em}

% ---------------------------------------------------------------
% Document-level metadata  (edit here for each new experiment)
% ---------------------------------------------------------------
\newcommand{\CourseCode}{EEE~2188}
\newcommand{\CourseTitle}{Electronics Sessional}
\newcommand{\ExperimentNo}{01}
\newcommand{\ExperimentName}{Familiarization of Different Equipment Involved with Electronics Sessional}

% ---------------------------------------------------------------
% Helper: lab equipment image
%   \labimage{path}{caption text}{label-key}
% ---------------------------------------------------------------
\newcommand{\labimage}[3]{%
  \begin{figure}[H]
    \centering
    \IfFileExists{#1}{%
      \includegraphics[width=0.55\textwidth,
                       height=5.5cm,
                       keepaspectratio]{#1}%
    }{%
      \fbox{\parbox[c][5.5cm][c]{0.55\textwidth}{\centering Missing image\\\detokenize{#1}}}%
    }%
    \caption{#2}
    \label{fig:#3}
  \end{figure}%
}

% Helper: wide theory diagram
\newcommand{\theoryimage}[3]{%
  \begin{figure}[H]
    \centering
    \IfFileExists{#1}{%
      \includegraphics[width=0.88\textwidth]{#1}%
    }{%
      \fbox{\parbox[c][5cm][c]{0.88\textwidth}{\centering Missing image\\\detokenize{#1}}}%
    }%
    \caption{#2}
    \label{fig:#3}
  \end{figure}%
}

% =============================================================
\begin{document}
% =============================================================

% ---------------------------------------------------------------
% Chapter 1: Experiment 01
% ---------------------------------------------------------------
\setcounter{page}{1}
\chapter{Familiarization of Different Equipment Involved with Electronics Sessional}
\markright{Objectives}

% =============================================================
\section{Objectives}
\markright{Objectives}
% =============================================================

\begin{enumerate}
  \item To learn about the commonly used equipment in the Electronics Sessional.
  \item To observe the various parameters of an electric circuit.
  \item To identify the different instruments ordinarily employed for electrical measurement.
  \item To understand the purpose and correct method of use of each instrument and component.
\end{enumerate}

% =============================================================
\section{Theory}
\markright{Theory}
% =============================================================

\subsection*{Background and Historical Context}

The discipline of electronics traces its origin to the late nineteenth century, when
physicists first began studying the behavior of electric charge in vacuum. The invention
of the thermionic (vacuum) diode by Sir John Ambrose Fleming in 1904, and the triode by
Lee de Forest in 1906, marked the dawn of active electronic devices capable of amplifying
signals \cite{horowitz2015}. For the following four decades, vacuum tubes powered radio
receivers, audio amplifiers, radar systems, and the earliest electronic computers. The
ENIAC computer (1945), for instance, contained 17{,}468 vacuum tubes and consumed roughly
150 kW of electrical power.

The pivotal transformation came in 1947, when William Shockley, John Bardeen, and Walter
Brattain at Bell Laboratories demonstrated the first transistor---a compact, solid-state
device built from germanium semiconductor material that could amplify current without a
vacuum envelope \cite{razavi2013}. This invention earned the trio the 1956 Nobel Prize in
Physics and laid the foundation for every modern electronic system. Silicon gradually
replaced germanium through the 1950s, and in 1958--1959, Jack Kilby (Texas Instruments)
and Robert Noyce (Fairchild Semiconductor) independently developed the integrated circuit
(IC), placing multiple transistors and passive components on a single semiconductor die.
Gordon Moore's 1965 observation---that the number of transistors on an IC doubles
approximately every two years---has held remarkably true for over five decades, enabling
the billions of transistors found on contemporary microprocessors \cite{sedra2020}.

Today, electronics is the enabling technology behind communication networks, medical
devices, power systems, transportation, and computing. The components studied in this
experiment---resistors, capacitors, diodes, transistors, and operational amplifiers---are
the fundamental building blocks from which all of these systems are constructed.

\subsection*{Fundamental Laws Governing Electronic Circuits}

All electronic circuits obey a small set of universal physical laws. \textbf{Ohm's Law}
states that the voltage $V$ across a resistive element is proportional to the current $I$
flowing through it: $V = IR$, where $R$ is the resistance in ohms. In a more general
sense, any two-terminal element is characterized by its voltage--current (V--I)
relationship, which may be linear (resistor), nonlinear (diode, transistor), or
energy-storing (capacitor, inductor) \cite{floyd2012}.

\textbf{Kirchhoff's Current Law (KCL)} states that the algebraic sum of all currents
entering any node of a circuit is zero (conservation of charge). \textbf{Kirchhoff's
Voltage Law (KVL)} states that the algebraic sum of all voltages around any closed loop
is zero (conservation of energy). These two laws, together with the V--I relationships of
the individual components, form a complete system of equations that uniquely determines
the voltages and currents at every point in a circuit \cite{floyd2012}.

Figure~\ref{fig:vi_char} illustrates the contrast between the linear V--I relationship
of a resistor and the strongly nonlinear exponential characteristic of a p--n junction
diode, described by the \textit{Shockley equation}:
\[
  I_D = I_S \left( e^{V_D / nV_T} - 1 \right)
\]
where $I_S$ is the reverse saturation current, $n$ is the ideality factor
($1 \leq n \leq 2$), and $V_T = kT/q \approx 26\,\text{mV}$ at room temperature is the
thermal voltage.

\theoryimage{images/vi_characteristics.pdf}%
  {V--I characteristics: linear resistor (left) and p--n junction diode
   showing the exponential forward-bias relationship and negligible reverse
   leakage current (right).}{vi_char}

\subsection*{Active vs.\ Passive Components}

Electronic components are broadly classified into \textbf{passive} and
\textbf{active} types. Passive components---resistors, capacitors, inductors, and
diodes---cannot supply energy to a circuit; they can only dissipate, store, or redirect
it. Active components---transistors (BJT and MOSFET) and operational amplifiers---can
control or amplify signals by drawing power from an external supply, delivering more
energy to the output than they receive at the input \cite{sedra2020}. This distinction is
fundamental: a passive filter can shape a signal but attenuates it; an active filter can
both shape and amplify it.

\subsection*{Measurement Instruments}

Understanding and verifying circuit behavior requires measurement. The two principal
instruments used in this sessional are the \textbf{multimeter} and the
\textbf{oscilloscope}. The multimeter provides accurate DC and AC scalar measurements
(voltage, current, resistance, continuity), while the oscilloscope displays the full
time-domain waveform of a signal, revealing its shape, frequency, amplitude, phase, and
transient behavior \cite{wiki:oscilloscope}. Figure~\ref{fig:scope_wave} shows a
sinusoidal waveform as displayed on a digital storage oscilloscope, with key measurement
parameters indicated. The oscilloscope is indispensable wherever signals change over time
faster than the human eye or a multimeter can track.

\theoryimage{images/oscilloscope_waveform.pdf}%
  {Sinusoidal waveform displayed on a digital storage oscilloscope,
   illustrating period $T$, peak voltage, and baseline.}{scope_wave}

% =============================================================
\section{Major Components and Equipment}
\markright{Major Components and Equipment}
% =============================================================

% -----------------------------------------------------------
\subsection{Resistor}
\markright{Resistor}
% -----------------------------------------------------------

\subsubsection*{Definition}
A resistor is a passive, two-terminal electronic component that implements a defined
electrical resistance as a circuit element \cite{wiki:resistor}. It obeys Ohm's Law
($V = IR$), producing a voltage drop proportional to the current through it, and converts
electrical energy irreversibly into heat. Resistance is measured in ohms ($\Omega$). For
through-hole resistors, the resistance value and tolerance are encoded in a system of
colored bands printed around the component body (see Figure~\ref{fig:resistor_lab}): each
color represents a digit (0--9), a multiplier, or a tolerance percentage. Common band
systems use 4, 5, or 6 bands; the mnemonic ``Bad Beer Rots Our Young Guts But Vodka Goes
Well'' (Black, Brown, Red, Orange, Yellow, Green, Blue, Violet, Grey, White) encodes the
digit order 0--9 \cite{electronics_notes:resistor}.

\begin{figure}[H]
  \centering
  \begin{subfigure}[b]{0.44\textwidth}
    \centering
    \includegraphics[width=\textwidth, height=5cm, keepaspectratio]{images/resistors.jpg}
    \caption{Axial-lead resistors}
  \end{subfigure}
  \hfill
  \begin{subfigure}[b]{0.52\textwidth}
    \centering
    \includegraphics[width=\textwidth, height=5cm, keepaspectratio]{images/resistor_colorcode.png}
    \caption{Color-band coding chart}
  \end{subfigure}
  \caption{Through-hole resistors and the standard color-band coding system used to
           identify resistance values and tolerances.}
  \label{fig:resistor_lab}
\end{figure}

\subsubsection*{Specifications}
Key resistor specifications include: resistance value (m$\Omega$ to G$\Omega$), power
rating ($\frac{1}{8}$\,W to several hundred watts for industrial types), tolerance
($\pm0.1\%$ for precision metal-foil to $\pm20\%$ for general-purpose carbon composition),
and temperature coefficient of resistance (TCR)---how resistance drifts with temperature,
expressed in ppm/$^\circ$C \cite{electronics_notes:resistor}. Metal-film resistors achieve
TCR values of $\pm25$\,ppm/$^\circ$C, carbon-film types are typically
$\pm100$ to $\pm500$\,ppm/$^\circ$C, and carbon-composition resistors may drift
$\pm1000$\,ppm/$^\circ$C.

\subsubsection*{Purpose of Use}
The fundamental purpose of a resistor is to limit or regulate the flow of electric
current. Beyond this, resistors divide voltages, set the bias (operating point) of
transistors, establish the gain of amplifier stages, terminate transmission lines to
prevent reflections, and provide pull-up or pull-down functions on digital logic lines.

\subsubsection*{Applications}
Resistors are found in every electronic system without exception. They appear in
sensor conditioning networks, transistor biasing circuits, voltage dividers in ADC
front-ends, LED current-limiting circuits, audio equalizers, feedback networks around
operational amplifiers, and as load banks for power supply testing \cite{boylestad2015}.

% -----------------------------------------------------------
\subsection{Multimeter}
\markright{Multimeter}
% -----------------------------------------------------------

\subsubsection*{Definition}
A multimeter---also called a volt-ohm meter (VOM)---is a multi-function electronic
measuring instrument that combines the capabilities of a voltmeter, an ammeter, and an
ohmmeter in a single portable unit \cite{wiki:multimeter}. Modern digital multimeters
(DMMs) convert measured quantities to digital form and display them on an LCD. Older
analog multimeters deflect a pointer over a calibrated scale using a moving-coil
D'Arsonval galvanometer mechanism.

\labimage{images/multimeter.jpg}%
  {A digital multimeter (DMM): the central rotary dial selects the
   measurement function and range; probes connect via colour-coded banana sockets.}%
  {multimeter}

\subsubsection*{Specifications}
A typical handheld DMM measures DC voltage (200\,mV to 600\,V), AC voltage (up to
600\,V), DC current (200\,$\mu$A to 10\,A), and resistance (200\,$\Omega$ to
200\,M$\Omega$) \cite{electronicsclub:multimeter}. Input impedance is standardized at
10\,M$\Omega$ to minimize circuit loading. DC voltage accuracy for a quality DMM is
typically $\pm(0.5\%\,\mathrm{rdg} + 1\,\mathrm{digit})$; laboratory bench instruments
achieve better than $\pm0.01\%$ of reading. Many DMMs also measure capacitance,
frequency, diode forward voltage, and transistor $h_{FE}$.

\subsubsection*{Purpose of Use}
The multimeter verifies operating voltages, checks component values, tests continuity
(detecting broken or short-circuit connections), measures the forward voltage of diodes,
and confirms the polarity of electrolytic capacitors. As a voltmeter it is connected
in parallel; as an ammeter it must be inserted in series via the dedicated current
terminal.

\subsubsection*{Applications}
Multimeters are indispensable in electronics laboratories for circuit debugging and
verification, in field service for diagnosing household wiring, automotive electrical
faults, and industrial control panels \cite{wiki:multimeter}. Clamp-meter variants
extend the same measurement philosophy to high AC currents without breaking the circuit.

% -----------------------------------------------------------
\subsection{Project Board / Breadboard}
\markright{Project Board / Breadboard}
% -----------------------------------------------------------

\subsubsection*{Definition}
A breadboard (or solderless breadboard) is a reusable construction base for prototyping
electronic circuits without soldering \cite{wiki:breadboard}. Beneath the plastic
surface, rows of spring-contact strips connect groups of holes electrically, allowing
component leads and jumper wires to be inserted and removed at will.

\labimage{images/breadboard.jpg}%
  {A full-size solderless breadboard (830 tie-points) with power rails
   (top and bottom) and component strips in the center section.}%
  {breadboard}

\subsubsection*{Specifications}
Holes are spaced at 0.1\,inch (2.54\,mm) pitch to match DIP integrated circuit packages.
A standard full-size breadboard provides 830 tie-points. Contact resistance is a few
ohms; inter-strip parasitic capacitance is a few picofarads, limiting reliable operation
to below roughly 10\,MHz \cite{analog:breadboard}.

\subsubsection*{Purpose of Use}
The breadboard enables rapid assembly, testing, and modification of circuits during the
design and verification phase. No tools are required: components are inserted by hand,
and reconfiguring a circuit is a matter of moving wires rather than removing solder.

\subsubsection*{Applications}
Breadboards are the standard prototyping platform in electronics education and in early
product development. They are used for assembling analog filters, digital logic circuits,
microcontroller (Arduino, Raspberry Pi) projects, and sensor interface circuits
\cite{analog:breadboard}. Once validated, designs move to veroboard or custom PCB.

% -----------------------------------------------------------
\subsection{Capacitor}
\markright{Capacitor}
% -----------------------------------------------------------

\subsubsection*{Definition}
A capacitor is a passive two-terminal component that stores and releases electrical energy
in an electric field \cite{sedra2020}. It consists of two conductive plates separated by
an insulating dielectric. The stored charge $Q$ relates to the applied voltage $V$ by
the capacitance $C$ (in farads, F):
\[
  C = \frac{Q}{V} = \frac{\varepsilon_r \varepsilon_0 A}{d}
\]
where $\varepsilon_r$ is the relative permittivity of the dielectric, $A$ is the plate
area, and $d$ is the plate separation (see Figure~\ref{fig:cap_structure}).

\theoryimage{images/capacitor_structure.pdf}%
  {Parallel-plate capacitor structure: two conductors separated by a dielectric
   store energy in the electric field between them.}{cap_structure}

\begin{figure}[H]
  \centering
  \begin{subfigure}[b]{0.44\textwidth}
    \centering
    \includegraphics[width=\textwidth, height=5cm, keepaspectratio]{images/capacitors.jpg}
    \caption{Electrolytic capacitors (polarized, high capacitance)}
  \end{subfigure}
  \hfill
  \begin{subfigure}[b]{0.44\textwidth}
    \centering
    \includegraphics[width=\textwidth, height=5cm, keepaspectratio]{images/ceramic_cap.jpg}
    \caption{Ceramic disc capacitors (non-polarized)}
  \end{subfigure}
  \caption{Common capacitor types used in electronics laboratory work. The lab
           contains 2200\,$\mu$F/25\,V and 10\,$\mu$F/16\,V electrolytic units alongside
           ceramic disc and polyester film types.}
  \label{fig:capacitors_lab}
\end{figure}

\subsubsection*{Specifications}
Capacitance values span from femtofarads (ceramic chip, high-frequency bypass) to tens
of thousands of microfarads (aluminum electrolytic, power supply filtering). Critical
specifications include capacitance, voltage rating (maximum safe DC voltage), tolerance,
equivalent series resistance (ESR), leakage current, and temperature coefficient
\cite{boylestad2015}. Electrolytic types are polarized (must be connected with correct
polarity); ceramic, film, and tantalum types are available in non-polarized forms.

\subsubsection*{Purpose of Use}
Capacitors block DC while passing AC (coupling), smooth power supply ripple (filtering),
store charge for pulsed loads (energy storage), set oscillation frequencies in RC timing
circuits, and form resonant circuits with inductors for frequency selection.

\subsubsection*{Applications}
Capacitors are among the most commonly placed components in any circuit. Power supply
filter capacitors reduce ripple; decoupling capacitors suppress high-frequency noise on
IC power pins; coupling capacitors transfer audio or RF signals between stages; and
timing capacitors set delays and oscillation frequencies in astable circuits
\cite{horowitz2015}.

% -----------------------------------------------------------
\subsection{Diode}
\markright{Diode}
% -----------------------------------------------------------

\subsubsection*{Definition}
A diode is a two-terminal semiconductor device formed by a p--n junction that allows
current to flow primarily in one direction \cite{boylestad2015}. The terminal attached
to the p-type material is the anode; the terminal at the n-type side is the cathode. In
forward bias (anode positive), the junction barrier is lowered and current flows; in
reverse bias, the barrier is increased and only a negligible leakage current flows until
reverse breakdown is reached. The internal structure of the p--n junction is illustrated
in Figure~\ref{fig:pnjunction}.

\theoryimage{images/pn_junction.pdf}%
  {Internal structure of a p--n junction diode showing the p-type region
   (holes), depletion region, n-type region (electrons), and conventional
   current direction under forward bias.}{pnjunction}

\labimage{images/diode_1n4007.jpg}%
  {1N4007 silicon rectifier diodes (DO-41 package) on a tape reel: rated
   1\,A average forward current and 1000\,V peak inverse voltage.}%
  {diode_lab}

\subsubsection*{Specifications}
Key specifications are: forward voltage drop $V_F$ ($\approx 0.6$--$0.7$\,V for silicon,
$\approx 0.2$--$0.4$\,V for Schottky), maximum average forward current $I_F$, peak
inverse voltage (PIV), and reverse recovery time $t_{rr}$ (important for high-frequency
switching) \cite{sedra2020}. The 1N4007---the most common general-purpose rectifier---has
$I_F = 1$\,A, PIV $= 1000$\,V, and $V_F \approx 1.1$\,V at rated current.

\subsubsection*{Purpose of Use}
The primary function of a diode is rectification: converting bidirectional AC to
unidirectional DC. Additional uses include voltage clamping (protecting circuits from
transient spikes), free-wheeling (absorbing inductive kickback), signal demodulation,
and logic level shifting.

\subsubsection*{Applications}
Diodes are found in AC--DC power supply rectifier bridges, battery charger protection
circuits (preventing reverse discharge), AM radio demodulators, inductive load protection
(flyback diodes across relay coils), and overvoltage protection at IC input pins
\cite{floyd2012}.

% -----------------------------------------------------------
\subsection{Zener Diode}
\markright{Zener Diode}
% -----------------------------------------------------------

\subsubsection*{Definition}
A Zener diode is a specially doped silicon p--n junction designed to operate
reproducibly in reverse breakdown at a controlled voltage, called the \textit{Zener
voltage} $V_Z$ \cite{boylestad2015}. At voltages below $V_Z$ the device blocks current
like a normal diode; once $V_Z$ is reached, current flows freely in reverse while the
voltage across the device remains nearly constant. This voltage-clamping property makes
the Zener diode the simplest voltage reference element.

\labimage{images/zener.jpg}%
  {Zener diode in DO-41 glass package; the band marks the cathode
   (negative terminal in the reference voltage context).}{zener_lab}

\subsubsection*{Specifications}
The Zener voltage $V_Z$ is available in standard values from about 2.4\,V to 200\,V.
Other key parameters are power dissipation rating $P_Z$ (commonly 400\,mW or 1\,W),
dynamic impedance $Z_Z$ (a few ohms---lower is better for regulation), and temperature
coefficient of $V_Z$ \cite{sedra2020}. Zener diodes with $V_Z \approx 5.1$\,V have a
near-zero temperature coefficient, making them the preferred choice for stable voltage
references.

\subsubsection*{Purpose of Use}
The Zener diode maintains a stable reference voltage across a load despite variations in
the input supply or load current. Placed in reverse bias with a series resistor, it
clamps the output to $V_Z$ as long as the Zener current remains within its operating
range.

\subsubsection*{Applications}
Zener diodes provide voltage references in linear power supplies, protect inputs from
transient overvoltage, set threshold voltages in comparator circuits, stabilize transistor
bias voltages, and regulate voltages in battery-powered devices \cite{boylestad2015}.

% -----------------------------------------------------------
\subsection{Light-Emitting Diode (LED)}
\markright{Light-Emitting Diode (LED)}
% -----------------------------------------------------------

\subsubsection*{Definition}
A light-emitting diode (LED) is a forward-biased p--n junction that converts electrical
energy directly into light through a process called \textit{electroluminescence}
\cite{sedra2020}. When electrons and holes recombine across the junction, they release
energy in the form of photons. The wavelength---and thus color---of the emitted light
depends on the semiconductor material's bandgap energy. Red and infrared LEDs use
AlGaAs or GaAsP; green and yellow use GaP or AlInGaP; high-brightness blue and white
LEDs (commercially available since the 1990s, earning Akasaki, Amano, and Nakamura the
2014 Nobel Prize in Physics) use InGaN \cite{razavi2013}.

\labimage{images/leds_photo.jpg}%
  {Standard 5\,mm through-hole LEDs in red, yellow, and green: the longer
   lead is the anode; the flat edge on the base identifies the cathode.}{led_lab}

\subsubsection*{Specifications}
Typical 5\,mm indicator LEDs have: $V_F = 1.8$--$2.2$\,V (red/yellow) or
$3.0$--$3.5$\,V (blue/white), maximum $I_F = 20$\,mA, luminous intensity of
10--1000\,mcd, and viewing angle of 15$^\circ$--60$^\circ$ \cite{horowitz2015}.
High-power LEDs rated at 350\,mA to several amperes are used in lighting applications.

\subsubsection*{Purpose of Use}
In the laboratory, LEDs are used as power and signal indicator lights. They must always
be paired with a current-limiting resistor calculated as
$R = (V_\text{supply} - V_F) / I_F$.

\subsubsection*{Applications}
LEDs appear in indicator panels, numeric and alphanumeric displays, LCD backlighting,
automotive headlamps and tail lights, traffic signals, optical fiber transmitters,
infrared remote controls, and solid-state general lighting \cite{razavi2013}.

% -----------------------------------------------------------
\subsection{Veroboard}
\markright{Veroboard}
% -----------------------------------------------------------

\subsubsection*{Definition}
A veroboard (stripboard) is a semi-permanent prototyping board made of copper-clad
fiberglass (FR4) laminate \cite{wiki:breadboard}. Unlike a breadboard, component leads
are soldered to the copper tracks, producing a robust and reliable circuit suitable for
regular use. Horizontal copper strips run between rows of holes; where tracks must be
broken, a track-cutting tool removes a small section of copper.

\labimage{images/veroboard.jpg}%
  {Veroboard (stripboard) showing continuous copper tracks on the
   underside and the grid of plated-through holes on the component side.}{veroboard_lab}

\subsubsection*{Specifications}
Standard veroboard uses 2.54\,mm (0.1-inch) hole pitch, 1.6\,mm FR4 substrate, and
35\,$\mu$m (1\,oz) copper tracks. Copper strips carry several amperes continuously
\cite{analog:breadboard}.

\subsubsection*{Purpose of Use}
Veroboard bridges the gap between breadboard prototyping and PCB fabrication. It is
chosen when a circuit must be portable and reliable but producing a custom PCB is
unnecessary or uneconomical.

\subsubsection*{Applications}
Veroboard is used for audio amplifier assemblies, sensor interface boards, hobbyist
project builds, educational laboratory assignments requiring a permanent record, and
small-batch production of simple circuits \cite{boylestad2015}.

% -----------------------------------------------------------
\subsection{Oscilloscope}
\markright{Oscilloscope}
% -----------------------------------------------------------

\subsubsection*{Definition}
An oscilloscope is an electronic test instrument that displays the variation of an
electrical signal (voltage) as a function of time \cite{wiki:oscilloscope}. It was
first developed in the 1890s as the cathode-ray oscilloscope (CRO), which deflected an
electron beam across a phosphorescent screen. Modern instruments are digital storage
oscilloscopes (DSOs): the analog input is sampled and digitized by an analog-to-digital
converter (ADC), stored in memory, and rendered on an LCD. DSOs can store waveforms for
later analysis, perform automatic measurements, and transfer data to a PC via USB or
LAN.

This laboratory is equipped with three oscilloscope models:

\begin{figure}[H]
  \centering
  \begin{subfigure}[b]{0.31\textwidth}
    \centering
    \includegraphics[width=\textwidth, height=4.5cm, keepaspectratio]{images/protek5100.jpg}
    \caption{Protek 5100 DSO\\(100\,MHz, 500\,MSa/s)}
  \end{subfigure}
  \hfill
  \begin{subfigure}[b]{0.31\textwidth}
    \centering
    \includegraphics[width=\textwidth, height=4.5cm, keepaspectratio]{images/rs_rtc1002.jpg}
    \caption{Rohde \& Schwarz RTC1002\\(300\,MHz, 2\,GSa/s)}
  \end{subfigure}
  \hfill
  \begin{subfigure}[b]{0.31\textwidth}
    \centering
    \includegraphics[width=\textwidth, height=4.5cm, keepaspectratio]{images/analog_oscilloscope.jpg}
    \caption{Analog oscilloscope\\(CRT display)}
  \end{subfigure}
  \caption{Oscilloscopes available in the EEE 2188 laboratory: (a) Protek 5100 DSO,
           (b) Rohde \& Schwarz RTC1002, and (c) analog CRT oscilloscope.}
  \label{fig:oscilloscopes}
\end{figure}

\subsubsection*{Specifications}
The most critical oscilloscope specification is \textit{bandwidth}---the $-3$\,dB
frequency at which a sinusoidal input is attenuated to 70.7\% of its true amplitude
\cite{tek:oscilloscope}. The Protek 5100 in this lab provides 100\,MHz bandwidth and
500\,MSa/s sampling rate; the Rohde \& Schwarz RTC1002 offers 300\,MHz and 2\,GSa/s.
As a rule of thumb, the oscilloscope bandwidth should be at least five times the highest
signal frequency of interest \cite{keysight:bandwidth}. Other key specifications include
vertical sensitivity (typically 1\,mV/div to 10\,V/div), time base
(ns/div to s/div), memory depth, number of channels, and trigger modes (edge, pulse,
video, serial bus).

\subsubsection*{Purpose of Use}
The oscilloscope reveals time-domain waveform information that a multimeter cannot
capture: signal shape, frequency, amplitude, phase, rise time, duty cycle, harmonic
content, and transient events. Two-channel instruments allow the simultaneous display of
input and output signals to measure gain and phase shift.

\subsubsection*{Applications}
Oscilloscopes are used in electronics design, manufacturing test, telecommunications
(signal integrity analysis), power electronics (switching transient measurement in
inverters), automotive diagnostics (CAN/LIN bus decoding), and biomedical engineering
(ECG, EEG signal acquisition) \cite{wiki:oscilloscope}.

% -----------------------------------------------------------
\subsection{Function Generator}
\markright{Function Generator}
% -----------------------------------------------------------

\subsubsection*{Definition}
A function generator is an electronic instrument that produces periodic waveforms of
controlled frequency, amplitude, and shape \cite{horowitz2015}. Common output waveforms
include sine, square, triangle, sawtooth, and arbitrary user-defined shapes. The
instrument in this laboratory is the \textbf{Protek 9205C Sweep Function Generator},
which also supports frequency sweep (varying the output frequency over a defined range),
useful for testing frequency-dependent circuits such as filters and amplifiers.

\labimage{images/protek9205c.jpg}%
  {Protek 9205C Sweep Function Generator used in this laboratory;
   front-panel controls set frequency, amplitude, waveform type, and
   sweep range.}{funcgen_lab}

\subsubsection*{Specifications}
The Protek 9205C covers output frequencies from below 1\,Hz up to approximately
5\,MHz across its sine, square, and triangle outputs, with output amplitude adjustable
from millivolts to several volts peak-to-peak into 50\,$\Omega$ loads. Frequency
accuracy is governed by a crystal-controlled timebase.

\subsubsection*{Purpose of Use}
Function generators supply the test signals needed to characterize amplifier gain and
bandwidth, filter frequency response, oscillator start-up behavior, and digital logic
timing margins. The sweep function automates frequency-response measurements.

\subsubsection*{Applications}
Function generators are used in amplifier characterization, filter testing, transducer
excitation, lock-in amplifier reference supply, and educational circuit experiments
across all frequency ranges \cite{horowitz2015}.

% -----------------------------------------------------------
\subsection{Operational Amplifier (Op-Amp)}
\markright{Operational Amplifier (Op-Amp)}
% -----------------------------------------------------------

\subsubsection*{Definition}
An operational amplifier (op-amp) is a high-gain, DC-coupled differential voltage
amplifier realized as an integrated circuit \cite{wiki:opamp}. It was originally
developed in the 1940s for use in analog computers, where it performed mathematical
\textit{operations} (summation, integration, differentiation) on voltages---hence the
name. The op-amp has two differential inputs (inverting~$-$ and non-inverting~$+$) and
one output; virtually all practical applications surround it with a negative feedback
network to set a precise closed-loop behavior \cite{horowitz2015}. The most iconic
general-purpose op-amp, the LM741, has been in production since 1968 and remains a
standard teaching device.

\labimage{images/lm741.jpg}%
  {LM741CN operational amplifier in DIP-8 package. Pin 2: inverting input;
   Pin 3: non-inverting input; Pin 6: output; Pins 7/4: $\pm$V supply.}{opamp_lab}

The inverting amplifier configuration---the most commonly used op-amp circuit---is shown
in Figure~\ref{fig:opamp_inv}. Negative feedback through $R_f$ forces the inverting
input to a \textit{virtual ground}, making the closed-loop gain precisely
$A_v = -R_f / R_\text{in}$, independent of the op-amp's open-loop gain.

\theoryimage{images/opamp_inverting.pdf}%
  {Op-amp inverting amplifier: the output is $180^\circ$ out of phase with
   the input, and the closed-loop voltage gain is $A_v = -R_f/R_\text{in}$.}{opamp_inv}

\subsubsection*{Specifications}
For an ideal op-amp: $A_{OL} = \infty$, $Z_{in} = \infty$, $Z_{out} = 0$, bandwidth
$= \infty$ \cite{electronics_tuts:opamp}. The practical LM741 has: open-loop gain
$\approx 200{,}000$, input impedance $\approx 2$\,M$\Omega$, output impedance
$\approx 75$\,$\Omega$, unity-gain bandwidth $\approx 1$\,MHz, slew rate
$0.5$\,V/$\mu$s, and CMRR of 70--90\,dB \cite{horowitz2015}.

\subsubsection*{Purpose of Use}
Op-amps amplify weak signals, implement mathematical operations, compare two voltages
(comparator mode), buffer high-impedance sources, and provide precise reference voltages.

\subsubsection*{Applications}
Op-amps are found in audio preamplifiers and equalizers, active filters, sensor signal
conditioning (strain gauges, thermocouples), analog-to-digital converter input stages,
signal generators, and industrial control loop compensators \cite{sedra2020}.

% -----------------------------------------------------------
\subsection{Bipolar Junction Transistor (BJT)}
\markright{Bipolar Junction Transistor (BJT)}
% -----------------------------------------------------------

\subsubsection*{Definition}
A bipolar junction transistor (BJT) is a three-terminal semiconductor device that
controls a large collector current $I_C$ with a small base current $I_B$ \cite{boylestad2015}.
It is built from two back-to-back p--n junctions sharing a thin \textit{base} region,
with \textit{emitter} and \textit{collector} regions on either side. NPN types (n-type
emitter and collector, p-type base) are more common. The BJT was invented at Bell Labs
in 1947 and was the primary amplifying device in all electronics for the following two
decades until the MOSFET gradually took over in digital applications.

\labimage{images/bjt_to92.jpg}%
  {Small-signal BJT in TO-92 plastic package. The three terminals---emitter,
   base, and collector---are identified by the flat face on the package
   and the manufacturer's datasheet.}{bjt_lab}

The BJT output ($I_C$--$V_{CE}$) characteristics, shown in Figure~\ref{fig:bjt_char},
display three distinct operating regions: the \textit{saturation} region (both junctions
forward-biased, used for switching), the \textit{active} region (base-emitter forward,
base-collector reverse, used for linear amplification), and the \textit{cutoff} region
(both reverse-biased, transistor off).

\theoryimage{images/bjt_characteristics.pdf}%
  {NPN BJT output characteristics ($I_C$ vs.\ $V_{CE}$) for five values
   of base current $I_B$. The saturation and active regions are indicated;
   $\beta = h_{FE} = \Delta I_C / \Delta I_B \approx 100$ in the active region.}{bjt_char}

\subsubsection*{Specifications}
Primary BJT parameters: DC current gain $h_{FE}$ (or $\beta$, ratio $I_C/I_B$, typically
50--500), collector--emitter breakdown voltage $V_{CEO}$, maximum collector current $I_C$,
maximum power dissipation $P_C$, and transition frequency $f_T$ (the frequency at which
current gain falls to unity) \cite{sedra2020}.

\subsubsection*{Purpose of Use}
In the \textit{active region} the BJT linearly amplifies current and voltage; in the
\textit{saturation/cutoff region} it operates as an electronic switch. The base--emitter
voltage $V_{BE}$ has a well-characterized temperature dependence of $-2$\,mV/$^\circ$C,
making BJTs useful as temperature sensors.

\subsubsection*{Applications}
BJTs are used in audio amplifiers, RF amplifiers, oscillators, relay and motor drivers,
TTL digital logic, temperature measurement circuits, and as the output stage of
operational amplifiers \cite{razavi2013}.

% -----------------------------------------------------------
\subsection{Metal-Oxide-Semiconductor Field-Effect Transistor (MOSFET)}
\markright{MOSFET}
% -----------------------------------------------------------

\subsubsection*{Definition}
A MOSFET is a voltage-controlled three-terminal semiconductor device that controls
current between its \textit{drain} (D) and \textit{source} (S) terminals via the
voltage applied to its electrically insulated \textit{gate} (G) \cite{wiki:mosfet}.
The gate is separated from the semiconductor by a thin silicon dioxide ($\text{SiO}_2$)
layer, making the gate current virtually zero and giving the MOSFET an extremely high
input impedance. MOSFETs are manufactured as n-channel (NMOS) and p-channel (PMOS)
types; complementary CMOS pairs underpin all modern digital logic. The circuit symbol
and terminal labels are shown in Figure~\ref{fig:mosfet_sym}.

\theoryimage{images/mosfet_symbol.pdf}%
  {N-channel enhancement-mode MOSFET circuit symbol: the arrow on the
   body indicates n-channel type; the dashed line shows the channel forms
   only when $V_{GS} &gt; V_{GS(\text{th})}$.}{mosfet_sym}

\begin{figure}[H]
  \centering
  \begin{subfigure}[b]{0.44\textwidth}
    \centering
    \includegraphics[width=\textwidth, height=5cm, keepaspectratio]{images/mosfet_to220.jpg}
    \caption{Power MOSFET in TO-220 package}
  \end{subfigure}
  \hfill
  \begin{subfigure}[b]{0.44\textwidth}
    \centering
    \includegraphics[width=\textwidth, height=5cm, keepaspectratio]{images/lm7805.jpg}
    \caption{LM7805CV voltage regulator (TO-220)}
  \end{subfigure}
  \caption{TO-220 packaged devices in this laboratory: (a) SSS30N06 power MOSFET
           (30\,A, 60\,V, Fairchild); (b) LM7805CV linear voltage regulator.}
  \label{fig:mosfet_lab}
\end{figure}

\subsubsection*{Specifications}
Key MOSFET parameters: gate threshold voltage $V_{GS(\text{th})}$ (the minimum
$V_{GS}$ to form a conducting channel), drain-to-source on-state resistance
$R_{DS(\text{on})}$ (milliohms for power devices, directly determining conduction loss),
maximum drain current $I_D$, breakdown voltage $V_{DSS}$, and transconductance $g_m$
\cite{razavi2013}. The SSS30N06 in this lab (Fairchild) is rated $I_D = 30$\,A,
$V_{DSS} = 60$\,V, $R_{DS(\text{on})} \approx 47$\,m$\Omega$.

\subsubsection*{Purpose of Use}
In the saturation region, the MOSFET is a voltage-controlled amplifier ($I_D \propto
(V_{GS}-V_{th})^2$). In the cut-off/linear (triode) region, it functions as a low-loss
electronic switch, conducting when $V_{GS} &gt; V_{th}$ and blocking when $V_{GS} &lt; V_{th}$.

\subsubsection*{Applications}
MOSFETs are the most produced electronic device in history. Billions of CMOS MOSFETs
constitute the logic gates in every processor, memory chip, and smartphone. Discrete
power MOSFETs switch kilowatts in motor drives, inverters, switch-mode power supplies,
and electric vehicle battery management systems \cite{wiki:mosfet}.

% -----------------------------------------------------------
\subsection{Voltage Regulator (LM7805CV)}
\markright{Voltage Regulator}
% -----------------------------------------------------------

\subsubsection*{Definition}
A linear voltage regulator is a three-terminal IC that maintains a fixed, stable DC
output voltage regardless of variations in the input supply or load current
\cite{sedra2020}. The LM7805CV is a member of the 78xx series, providing a fixed
$+5$\,V output. Internally it uses a Zener-diode reference, an error amplifier, and a
pass transistor connected as an emitter follower to regulate the output.

\labimage{images/lm7805.jpg}%
  {LM7805CV $+5$\,V linear voltage regulator in TO-220 package. Pins
   from left: input, ground, output.}{lm7805_lab}

\subsubsection*{Specifications}
The LM7805 accepts input voltages from 7\,V to 35\,V and delivers a regulated $+5$\,V
output at up to 1\,A continuous current ($1.5$\,A with heatsink). Line regulation is
typically 3\,mV (output change for full input swing); load regulation is about 15\,mV
\cite{boylestad2015}. A 0.33\,$\mu$F input decoupling capacitor and 0.1\,$\mu$F output
capacitor are recommended for stability.

\subsubsection*{Purpose of Use}
The LM7805 converts an unregulated DC supply (e.g., a 9\,V battery or 12\,V rectified
supply) into a stable $+5$\,V rail for powering digital logic circuits, microcontrollers,
and sensor modules.

\subsubsection*{Applications}
The 78xx family is used in laboratory bench power supply circuits, embedded system
prototypes, battery-powered sensor nodes, and educational PCB designs \cite{floyd2012}.

% -----------------------------------------------------------
\subsection{IC in DIP Package (AT28C64U / General-Purpose)}
\markright{IC in DIP Package}
% -----------------------------------------------------------

\subsubsection*{Definition}
An integrated circuit (IC) in a dual in-line package (DIP) is a semiconductor device
in which multiple transistors, resistors, and other elements are fabricated on a single
silicon die and enclosed in a rectangular plastic or ceramic body with two parallel rows
of pins on 2.54\,mm pitch \cite{sedra2020}. The DIP format is compatible with both
breadboard and veroboard prototyping. The 28-pin DIP device in this laboratory
(tentatively identified as an AT28C64U type EEPROM) is a representative example of how
digital memory and logic ICs are packaged for laboratory use.

\labimage{images/dip28_ic.jpg}%
  {A 28-pin DIP integrated circuit as used in the EEE 2188 laboratory.
   The notch at one end identifies pin 1; pins are numbered
   counter-clockwise from pin 1.}{dip28_lab}

\subsubsection*{Specifications}
DIP packages are available from 6 to 64 pins. The 28-pin DIP has an overall body width
of 0.6\,inch (15.24\,mm) and an overall pin-to-pin span of 0.1\,inch pitch. IC
specifications vary entirely by device type---digital, analog, memory, or mixed-signal.

\subsubsection*{Purpose of Use}
DIP ICs are inserted into breadboard or veroboard for prototyping digital logic, memory,
microcontroller, and analog IC circuits. The DIP format makes pin identification and
probing straightforward.

\subsubsection*{Applications}
DIP-packaged ICs include logic gates (74xx TTL/CMOS series), op-amps (LM741, LM324),
comparators (LM339), timer ICs (NE555), SRAM and EEPROM memory chips, and
microcontrollers (PIC, ATmega) \cite{sedra2020}.

% -----------------------------------------------------------
\subsection{Soldering Iron}
\markright{Soldering Iron}
% -----------------------------------------------------------

\subsubsection*{Definition}
A soldering iron is a hand-held tool with an electrically heated metal tip (the
\textit{bit}) used to melt solder and join component leads to printed circuit board
pads or wire terminals \cite{wiki:solderiron}. Solder is a low-melting-point metal
alloy that, when melted and solidified, forms a mechanically strong and electrically
conductive joint. Traditional tin--lead solder (60Sn/40Pb) melts at approximately
183\,$^\circ$C; modern lead-free alloys (e.g., Sn96.5/Ag3.0/Cu0.5, ``SAC305'') melt at
217--227\,$^\circ$C \cite{wiki:solderiron}.

\labimage{images/soldering_iron.jpg}%
  {A general-purpose soldering iron ($\approx$30--40\,W) of the type used
   in this laboratory; the replaceable tip (bit) is heated resistively to the
   set temperature.}{solderiron_lab}

\subsubsection*{Specifications}
The laboratory soldering iron is rated approximately 30--40\,W, operating at mains
voltage (220\,V AC). Tip temperature for general electronics work with lead-free solder
is set to 315--370\,$^\circ$C. Temperature-controlled soldering stations maintain the
set temperature to within a few degrees via a closed-loop thermocouple sensor, reducing
the risk of thermal damage to heat-sensitive components \cite{wiki:solderiron}.

\subsubsection*{Purpose of Use}
The soldering iron is used to create permanent solder joints on veroboard and PCBs,
to tin (pre-coat) wire ends and component leads, and to desolder components for rework.
A good solder joint is smooth, shiny, and conical; a defective ``cold joint'' is dull
and grainy, with unreliable electrical conductivity.

\subsubsection*{Applications}
Soldering irons are used in PCB assembly and repair, wire-to-terminal connections,
cable harness manufacturing, educational laboratory circuit assembly, and field-service
rework of electronic equipment \cite{wiki:solderiron}.

% =============================================================
\section{Discussion}
\markright{Discussion}
% =============================================================

In this experiment, thirteen categories of electronic components and measuring instruments
were studied: the resistor, multimeter, breadboard, capacitor, diode, Zener diode, LED,
veroboard, oscilloscope, function generator, operational amplifier, BJT, MOSFET, voltage
regulator, and soldering iron. Each component was examined physically, and its working
principle, key specifications, and laboratory applications were documented.

The distinction between passive components (resistors, capacitors, diodes) and active
devices (BJTs, MOSFETs, op-amps) is of fundamental importance: passive components
can only dissipate or store energy, while active devices can control or amplify signals
by drawing on an external power supply. Among the measurement instruments, the
oscilloscope proved especially versatile, capable of revealing waveform shape, frequency,
amplitude, and phase in a way that a multimeter alone cannot provide. The three
oscilloscope models in the lab---the Protek 5100 (100\,MHz), the Rohde \& Schwarz RTC1002
(300\,MHz), and the analog CRT unit---illustrate the progression from analog to digital
storage technology in oscilloscope design.

A consistent practical observation is the importance of component identification before
use: resistor color bands must be read correctly (or verified with the multimeter's
resistance function), electrolytic capacitors must be oriented with correct polarity,
and diode anode/cathode must be distinguished to avoid circuit damage. The LED requires
a current-limiting resistor in all configurations; without it, the device will be
destroyed in seconds.

% -- Students may extend the discussion here --

\vspace{2cm}
% (Continue below as needed)

% =============================================================
\section{Conclusion}
\markright{Conclusion}
% =============================================================

This experiment provided a comprehensive introduction to the components, instruments,
and tools that constitute the core equipment of the Electronics Sessional. By physically
handling and examining each item---identifying its terminals, reading its markings, and
connecting it to measurement instruments---the groundwork was established for all
subsequent experiments in EEE 2188.

The resistor, the most fundamental passive component, was shown to limit current
according to Ohm's Law; the capacitor stores charge and blocks DC; the diode and Zener
diode enforce one-directional current flow and voltage clamping respectively; and the
LED converts electrical energy to visible light. Among the active devices, the BJT and
MOSFET both act as amplifiers and switches, with the MOSFET's near-zero gate current
making it the dominant device in modern digital and power electronics. The operational
amplifier, configured with external feedback resistors, provides precise, stable
amplification suitable for a wide variety of signal processing tasks.

The breadboard and veroboard support circuit prototyping at different levels of
permanence, while the multimeter and oscilloscope provide the measurement capability
needed to verify and debug circuit behavior. The soldering iron completes the toolkit,
enabling permanent assembly of verified designs.

Familiarity with each of these items---their specifications, limitations, and correct
usage---is a prerequisite for the remaining experiments in this course and for
professional electronics engineering practice.

% -- Students may extend the conclusion here --

\vspace{1.5cm}

% =============================================================
% Bibliography
% =============================================================
\markright{References}
\bibliography{../../tex/references}

\end{document}
